\backhead{FURTHER EXERCISES}{khred}
\begin{exlist}
\item Using the data of the net-radiation example of this chapter, find $R_n$ if the sunshine falls to 4 h (same day and temperatures).\exans{10.01 MJ\,m$^{-2}$\,d$^{-1}$}
\item A soil has a thermal diffusivity of $3\times10^{-7}$ m$^2$\,s$^{-1}$. Find the daily damping depth and the depth at which the daily amplitude is one tenth of the surface amplitude.\exans{0.091 m; 0.21 m}
\item A canopy with $L=4$ stores 0.2 mm per unit leaf area. Of a 10 mm rain, how much is intercepted and how much reaches the soil? Of five rains of 2 mm, separated by drying?\exans{0.8 and 9.2 mm; 4.0 and 6.0 mm}
\item A snowpack of 120 mm water equivalent melts with $C_m=4$ mm\,$^\circ$C$^{-1}$\,d$^{-1}$ at 5 $^\circ$C. How long does it last? How much drains below a root zone that can store 60 mm more?\exans{6 d; 60 mm}
\item Find the capillary rise from a water table at 120 cm for a soil with $K_s=20$ cm\,d$^{-1}$ and $a=0.03$ cm$^{-1}$.\exans{5.62 mm\,d$^{-1}$}
\end{exlist}
\backhead{ESSAY QUESTIONS}{khred}
{\small\itshape Answer each question in about one page, with equations, a sketch and at least one number from this chapter.}\par\vspace{3pt}
\begin{exlist}
\item Explain the process map of Daisy and the four kinds of line. How does the map help to read a simulation?
\item Describe the chain from the weather to the potential evapotranspiration of the crop and the soil in Daisy.
\item Discuss the water stores above the soil (snow, canopy, litter, pond) and the order in which they are filled and emptied.
\item Explain how heat moves in the soil and why the soil temperature matters for water and nitrogen.
\item Discuss the role of biopores and drains in the water balance of a field.
\item Compare four equations of the reference evapotranspiration and discuss when each should be used.
\end{exlist}
